\chapter{最後の散乱とバリオン生成}

\section{外を見ることは過去を見ること}

距離が異なれば、我々は異なる時代の宇宙を見ることになる。近傍では、遡及時間はそれほど重要ではない。一年前や百年前の宇宙は、現在の宇宙とそれほど変わっていなかった。しかし十分遠方では、我々はより若い宇宙を見ることになる。より若い宇宙はより高密度であり、より高密度な宇宙はより高温である。

ある十分に大きな距離では、宇宙は太陽の表面とほぼ同じくらい高温であった。光子がここに届くまでに赤方偏移しているため、私たちにはそのようには見えない。赤方偏移因子がおよそ1000であれば、観測される温度は同じ因子だけ低くなる：
\[
T_{\mathrm{obs}} \sim \frac{T_{\mathrm{emit}}}{1+z},
\qquad
1+z \sim 10^3.
\]\footnote{この表示は00:02:20--00:02:39における口頭の赤方偏移の議論を簡潔に編集した転記である。}

そこが宇宙が不透明になった場所である。宇宙は主に水素で、少量のヘリウムを含んでいた。数千度を超える温度では、水素は電離している。したがってその表面より向こうには中性原子ではなくイオン、主に自由に運動する陽子と電子がある。そこでは光子もより高エネルギーで、基本的にX線である。原子が偶然形成されても、衝突または光子によってすぐに再電離される。

電離気体は電気伝導体であり、伝導体は放射をきれいに通さない。電荷は電場に応答し、前後に揺れ動き、放射を再放出する。したがってその距離より向こうの宇宙は私たちにとって不透明である。最後の散乱面は太陽の可視表面のようなものである：物質の壁ではなく、媒質がかなり突然に不透明から透明へ変わる場所である。

そこに存在する粒子のほとんどは光子であり、その割合はおよそ\(10^8\)倍である。同じことは今日にも当てはまるが、今日では光子のエネルギーが非常に低いため、現在のエネルギー密度の重要な部分ではない。温度が水素の電離温度を下回ると、宇宙は透明になり、私たちはそれを通して見ることができる。

直接の視界はそこで終わる。次の問題は、既知の原子物理学、核物理学、素粒子物理学がどこまでさらに過去へ遡れるかである。

\section{さらに過去へ遡る}

最後の散乱のすぐ内側では、光子に加えて、主な粒子は陽子と電子であり、いくつかのヘリウム原子核と少数の重陽子がある。反陽子は事実上なく、陽電子も事実上ない。しかしさらに過去へ遡ると、温度は最終的に2つの光子が衝突して電子と陽電子を生成するのに十分高くなる：
\[
\gamma + \gamma \leftrightarrow e^- + e^+.
\]

基本的なしきい値は、光子が少なくとも対の静止エネルギーを供給しなければならないことである：
\[
E \ge 2m_e c^2.
\]
そのしきい値以下では通常の光は役に立たない。それ以上では電子--陽電子対が生成され、逆反応も起こりうる。その温度にある高温の物質の鍋には、通常の物質で化学反応が果たす役割を粒子反応が果たしているものが含まれるであろう。

しきい値のすぐ上では、依然として光子が電子と陽電子より多い。電子の静止質量は高くつく。しかし温度がしきい値をはるかに超えて上昇すると、電子の静止エネルギーは運動エネルギーより重要でなくなり、密度は同程度になる。そのとき電子、陽電子、光子はすべてほぼ同じ密度を持ち、狂ったように衝突し合っている。

\begin{figure}[ht]
\centering
\includegraphics[width=0.62\textwidth]{lecture_06_figure_02.png}
\caption{電子と陽子の質量スケールを比較する対生成しきい値のスケッチ。}
\lectureframe{00:18:47}
\end{figure}

次の問題は陽子に何が起こるかである。陽子はすでに存在し、その時点では非常に密であるが、陽子--反陽子対を大量に生成するのに十分な熱エネルギーはまだない。理由は単純である：
\[
m_p \approx 2000\,m_e,
\qquad
E \ge 2m_p c^2
\]
はるかに高いしきい値である。\editorialnote{下の質量のラベルは上の\(2m_e c^2\)よりも鮮明でないが、口頭の比較は明示的に電子と陽子のしきい値の間で行われている。}

これはすでに恒星の内部よりもはるかに高温であり、そこでは反粒子は実質的な量では存在しない。

しかしそれさえも暫定的な記述にすぎない。温度が陽子と反陽子を自由に生成するのに十分高くなる頃には、陽子をクォークに粉砕するのにも十分高くなっている。したがって本当に高温のスープは、電子、陽電子、クォーク、反クォーク、光子、ニュートリノ、グルーオン、その他の種のスープとして記述する方がよい。黒板からの簡単な模式的反应は
\[
e^- + e^+ \to \gamma \to q + \bar q.
\]

\begin{classroomqa}
\lecturetimestamp{00:20:05}
\audiencequestion{それは中間子か、それともクォーク以前か？}
\lecturerresponse{中間子はクォークからできている。温度が陽子を粉砕するのに十分高くなれば、中間子も粉砕する。}
\end{classroomqa}

これまでの推論は時間を遡る方向だった。ここで向きを変え、高温で対称なスープが冷えていくときに何が起こるかを問おう。

\section{電気的中性と対称な高温宇宙}

十分に高い温度では、宇宙には粒子と反粒子が膨大な密度で存在する。電子と陽電子は対で現れ、クォークと反クォークは対で現れ、それぞれの粒子とその反粒子がほぼ等しい存在量であるというのが自然な描像である。

電気的中性には少し注意が必要である。今日の宇宙は電気的に中性である。有用な記憶術は、我々の実際の宇宙ではなく、三次元球面のような閉じて有界な宇宙を想像することである。まず二次元球面を模型として考える。電気力線は電荷以外では始まったり終わったりできない。正電荷は力線を放出し、負電荷はそれらを吸収する。閉じた球面上では、それらの線は単に無限遠へ逃げることはできない。それらは反対側で再び集束する。したがって記憶術として、閉じて有界な宇宙は正味の電荷を持ちえない。我々の実際の宇宙が文字通り閉じているかどうかは要点ではない。要点は、大規模な正電荷と負電荷が驚くほどよく相殺しているように見えることである。

電気的中性は、陽子の数が反陽子の数と同じであることを\emph{意味しない}。それは、正の全電荷が負の全電荷に等しいことを意味する。陽子、陽電子、その他あらゆる正に帯電した種は、電子およびその他あらゆる負に帯電した種と全体として釣り合う。

既知の物理学が信頼できる限り過去に遡っても、おそらくは少しそれを超えても、宇宙は粒子と反粒子の両方を豊富に維持できるほど高温だった。挙げられるどんな加速器実験にも、初期宇宙の何らかの段階に対応するものがあった。ほぼ対称な始まりが自然な推測である。すなわち、同数の電子と陽電子、同数のクォークと反クォーク、そして陽子と反陽子をバリオンと反バリオンの代役として用いるならば、
\[
N_p = N_{\bar p}.
\]

その推測は、物理学の法則が粒子と反粒子の間でほとんど対称であるため魅力的である。陽電子は電子と非常によく似た振る舞いをする。完全にではないが、ごく近い。その近似的対称性こそが問題である。

\section{映画を前へ進める}

宇宙が、粒子と反粒子が等しいかほぼ等しい存在量を持つ高温状態で始まり、その後冷えるとしよう。

温度が十分に高い間は、陽子と反陽子は高エネルギー衝突から生成され、光子へと対消滅して戻ることができる。しかしやがて温度は陽子のしきい値を下回る。すると陽子--反陽子対は依然として光子へ対消滅できるが、少し膨張した後にはそれらの光子はエネルギーを十分に失い、もはや陽子--反陽子対へ戻ることができなくなる。これが一方通行の道である。同じことが後に電子と陽電子について起こる。陽子のしきい値よりずっと低い温度だが、それでも恒星の温度よりはるかに高い。

膨張は微視的な反応率に比べて遅いので、これが起こっている間、宇宙はほとんど熱平衡に近い。これがここで断熱膨張によって意味されていることである。その状況では、より正確な言明はエントロピーがごくわずかしか変化しないということである。高温の粒子スープでは、エントロピーは本質的に粒子の数であり、1の程度の因子を除いてそうである。したがって、対消滅する陽子、反陽子、電子、陽電子は、おおよそ光子によって置き換えられる。

陽子と反陽子の数が厳密に等しければ、それらは単に対消滅して消えてしまうだろう。同じことが電子と陽電子にも起こるだろう。残るのは光子である。しかし陽子が反陽子よりごくわずかでも過剰であれば、対になった組は対消滅し、対にならない過剰分は残るだろう。具体的な例が要点を示す。陽子が百個、反陽子が五十個あれば、五十個の反陽子は五十個の陽子を見つけてそれらと対消滅し、五十個の陽子が残るだろう。電荷の中性は、次に電子が陽電子より同種の残存過剰を持つことを要求する。

これにより、現在の光子数から初期の非対称性を読み取ることができる。今日、光子の数は陽子または電子の数の約\(10^9\)から\(10^{10}\)倍である。
\[
\frac{N_\gamma}{N_p} \sim 10^9\text{--}10^{10}.
\]
したがって初期の高温スープでは、陽子--反陽子対十億組につき約一つの余分な陽子しかなかった。そのわずかな不均衡こそが、原子が今なお存在する理由のすべてである。

\begin{classroomqa}
\lecturetimestamp{00:42:25}
\audiencequestion{余剰を統計的に、冷却の際にごくわずかな確率で生成されたものとして考えてはどうだろうか。}
\lecturerresponse{最初に対称な状態から余分な陽子を得るには、反陽子も同時に生成することなく陽子を生成しなければならない。またそれとともに負の電荷、たとえば電子も生成しなければならない。既知の標準模型にはそのような低エネルギー過程は存在しない。これまでに検証されたエネルギーではバリオン数は保存されるので、通常の冷却は正味の余剰を生み出さない。}
\end{classroomqa}

したがって電荷は障害ではない。模式的な例は
\[
\gamma + \gamma \to p + e^-.
\]\footnote{この表示は00:44:16--00:46:35における口頭の例を編集上書き起こしたものである。その目的は、電荷が保存されつつバリオン数が破れうるという点を切り出すことにある。}
これは全体として電気的に中性であろう。しかし別の経験的規則に反する。検証された標準模型では、反陽子も同時に生成することなく陽子を生成することはできない。当面の目的にとって関連する保存量はバリオン数であり、
\[
B=(N_p+N_n)-(N_{\bar p}+N_{\bar n}).
\]\footnote{この簡潔な式は00:45:18--00:45:40で与えられた口頭の定義を定式化したものである。}
既知のエネルギーではその量は保存されるので、非対称性は冷却史の後期のよく検証された部分で作り出すことはできない。

\section{バリオン生成}

これがバリオン生成と呼ばれる主題である。宇宙が本当に等しい数のバリオンと反バリオンで始まったとしよう。その対称な状態を観測されたわずかな余剰へと傾けるために、物理法則は何を許さなければならないだろうか。

黒板ではこれは二つの記号式に還元され、
\[
N_p=N_{\bar p},
\qquad
N_p-N_{\bar p},
\]
次に第一の式を第二の式に結びつける条件の一覧に還元された。

\begin{figure}[ht]
\centering
\includegraphics[width=0.72\textwidth]{lecture_06_figure_04.png}
\caption{最初に対称な状態からバリオン非対称性を生成するための条件。}
\lectureframe{01:22:26}
\end{figure}

\editorialnote{このフレームでは番号付きリストの語は一部しか判読できない。以下の三項目は付随する口頭の議論に従っている。}

第一の条件はバリオン数の破れである。バリオン数が厳密な保存則であるならば、正味のバリオン数がゼロで始まった宇宙はそのままゼロにとどまる。

第二の条件は粒子と反粒子の間の偏りである。バリオン数の破れだけでは十分でない。もし自然が陽子を生成する過程とその反映である反陽子を生成する過程を等しい確率で許すならば、平均の不均衡は依然としてゼロである。生成率が異なることが必要である。

通常の統計的揺らぎもまたはるかに小さすぎる。光子を含めて粒子の総数がおよそ\(10^{90}\)であるならば、\(\sqrt{N}\)程度の揺らぎはおよそ\(10^{45}\)しか与えず、宇宙にあるおよそ\(10^{80}\)個の陽子を説明するにはまったく足りない。

第三の条件は、宇宙がかつて十分に非平衡であったことである。ここで議論はCPT定理を用いる。高温の気体が永遠に平衡にとどまっているならば、逆回しにした映画は順回しの映画と同じに見える。その状況では、粒子と反粒子の間の対称性は、電荷共役、パリティ、時間反転を組み合わせた厳密な対称性によって実効的に回復される。実際の歴史に真の偏りが凍結されるためには、実際の歴史が微視的な時間スケールで前向きと後ろ向きで異なって見えなければならない。十分に初期には、宇宙の膨張はそれを行えるほど速かった。

したがって三つの条件は次のとおりである。
\begin{enumerate}
\item バリオン数の破れ、
\item 粒子と反粒子の間の対称性の破れ、より正確にはCPの破れ、
\item 熱平衡から外れた段階。
\end{enumerate}

\begin{classroomqa}
\lecturetimestamp{01:09:08}
\audiencequestion{これらをCPTと対応させると、条件2は本当にCPであり、条件3は本当に時間の向きの問題なのか。}
\lecturerresponse{はい。条件2は本当はCPの破れとして述べられるべきである。条件3は、実際の歴史が非平衡であったという事実であり、別個の保存則でもなく、バリオン数を破ることと同じものでもない。}
\end{classroomqa}

これら三つの条件が揃えば、等しい数の陽子と反陽子で始まった宇宙でさえ、どちらか一方に傾けることができる。依然として不確かなのは、自然が本当にバリオン数を破るかどうか、そして破るならばどのように破るかである。

\section{何が分かっており、何が分かっていないか}

この話の未解決部分はバリオン数破れである。自然界で
\[
(N_p+N_n)-(N_{\bar p}+N_{\bar n}).
\]
を変える反応を直接見たことは一度もない。それは、関係する反応が実験室で直接試されたどの温度よりもはるかに高い温度を必要とするだけなのかもしれないし、低エネルギーでは途方もなく起こりにくいだけなのかもしれない。しかし現時点では、これは依然として疑問符のついた問題である。

だからこそ、この議論では陽子崩壊がそれほど重要なのである。もしバリオン数が厳密でないなら、陽子は原理的には不安定であるはずだ。しかし陽子は驚くほど安定である。ここで引用されている下限は
\[
\tau_p \gtrsim 10^{34}\ \text{years}.
\]
の程度である。これは、光検出器を備え付けられ深く地下に埋められた非常に大量の水を用いた巨大な実験から得られたものである。これはバリオン数破れを否定するものではない。そのような破れは低エネルギーでは並外れて抑制されていなければならない、ということを意味するだけである。

既知の物理に関する別の注意点もこの議論に入ってくる。\(100\)から\(250\,\mathrm{GeV}\)のあたりで、弱い相互作用は電磁相互作用にはるかに似たものになる。これは詳細な模型では重要かもしれないが、現在の議論における劇的な段階ではない。

暗黒物質は基本的な論理をあまり変えない。ほとんどの理論では、暗黒物質は通常の物質と非常に弱くしか相互作用しないので、この問題に関しては本質的に傍観者である。

\begin{classroomqa}
\lecturetimestamp{01:24:49}
\audiencequestion{重力はこのすべてにどのように影響するのか。}
\lecturerresponse{主に二つの仕方で影響する。第一に、宇宙膨張がなければ非平衡の段階は存在しない。第二に、ブラックホール蒸発は、バリオン数が永遠に厳密な保存則ではありえないと疑う強い理由を与える。それ以外では、粒子間の通常の重力相互作用はここで重要になるには弱すぎる。}
\end{classroomqa}

次に電荷についての追加の質問がある。同じブラックホールの議論は電荷の保存も脅かすのだろうか。答えは否である。帯電したブラックホールは外側へ延びる電気力線を伴っており、それらの力線が電荷を記録している。そのようなブラックホールは完全に蒸発しきることはなく、むしろ極限残骸に近づく。したがって、ブラックホールの議論が脅かすのはバリオン数であり、電荷ではない。

別の質問は、非常に高密度の初期のクォークスープを単にブラックホールと見なすべきかどうかを問うている。答えは否である。要点は静的な高密度の塊ではない。宇宙は膨張していた。ブラックホールの話はバリオン数の運命についての思考実験であり、初期宇宙を閉じ込められた領域として文字通り記述したものではない。

\begin{classroomqa}
\lecturetimestamp{01:31:14}
\audiencequestion{そのような反応を含めるように標準模型を修正する方法は分かるのだろうか。}
\lecturerresponse{困ったことは逆である。バリオン数破れを加えすぎて陽子をあまりに速く崩壊させるのは容易である。本当の問題は、望まれる初期宇宙の過程をなお許しながら、陽子を現状のように安定に保つことである。}
\end{classroomqa}

\section{エントロピーと後期の中性子に関する質問}

ここで簿記の話に戻ろう。断熱的に膨張する宇宙では、エントロピーは近似的に保存される。素粒子の熱いスープでは、エントロピーは本質的に粒子数であり、1の程度の因子を除いてそうである。だからこそ、この議論ではどこからともなく大量の光子が後期に生成される必要はない。対消滅は主にスープの中身を放射へと組み替えるのである。

問題の古い版では話が逆さまになっていた。光子に対してなぜこれほど少数の陽子が生き残ったのかを問うのではなく、逆に問われていた。なぜ陽子に比べてこれほど多くの光子が存在するのか、と。しかしより鋭い問いは、なぜ宇宙がゲームの後半でこれほど多くのエントロピーを生成したのか、ということではない。最初から粒子と反粒子の非対称性がなぜこれほどわずかだったのか、ということである。すなわち、10億個の陽子--反陽子対あたりおよそ1個だけ余分な陽子がある、という程度である。

自由中性子は実際に崩壊するが、その過程はバリオン数を\emph{破らない}：
\[
n \to p + e^- + \bar\nu.
\]
中性子は陽子と同じバリオン数を担っている。

自由中性子では利用可能な崩壊エネルギーが非常に小さい。中性子の質量が最終状態の質量よりわずかに大きいだけだからである。素粒子の基準ではこの崩壊がこれほど遅い理由の一部はそこにある。中性子が原子核に束縛されていると、束縛エネルギーがそのわずかな利用可能エネルギーを完全に取り除くことができ、そのとき束縛中性子は安定になりうる。

\begin{classroomqa}
\lecturetimestamp{01:43:03}
\audiencequestion{中性子は単に宇宙から消えているのか、それとも生成もされているのか。}
\lecturerresponse{第一の答えは、自由中性子は崩壊によって失われるということである。しかしそれだけがすべてではない。恒星内の核反応や中性子星の形成など、中性子生成の実際の源が存在する。したがって中性子の総数は単純に一方向に減少するだけではない。}
\end{classroomqa}
